% SGA 3, Expose VI_A, Loop 1 English reconstruction.
% Scope: Section 1 heading through the transition after Proposition 1.3.1;
% Exemples 1.3.2 is excluded.
% Authority: Polo--Gille Exp6A-13oct24.pdf, component-local pages 4--6
% (combined-reader pages 307--309; printed pages 296--298;
% current-rescribe indices 292--295).
% OCR and the Reinhold English Markdown are comparison witnesses only.

\section*{1. 局部有限型 \(A\)-群的局部性质}
\phantomsection\label{sec:sga3-VIA-1}
\addcontentsline{toc}{section}{1. 局部有限型
  \texorpdfstring{\(A\)}{A}-群的局部性质}

首先将说明, 若 \(G\) 局部有限型且在 \(A\) 上平坦,
则可通过基的一个有限平坦扩张,
``使 \(G\) 的每个闭点都严格有理''.

\medskip
\phantomsection\label{par:sga3-VIA-1.1}%
\noindent\textbf{1.1.}
从现在起, 除非明确另作说明,
都假设 \(G\) 是局部有限型的 \(A\)-群.
当 \(A\) 是域 \(k\) 时, 我们将在第 VII\textsubscript{B} 篇中
得到关于 \(G\) 的局部环的十分精确的结果.%
\footnote{\label{note:sga3-VIA-10}编者注: 例如, 它们总是完全交;
见第 VII\textsubscript{B} 篇, \SGARefLink{sga3:VIIA:5.5.1}{5.5.1}.
此外, 若 \(\operatorname{char}(k)=0\), 则 \(G\) 平滑
(第 VI\textsubscript{B} 篇, \SGARefLink{cor:sga3-VIB-1.6.1}{1.6.1};
另见第 VII\textsubscript{B} 篇, \SGARefLink{sga3:VIIA:3.3.1}{3.3.1}).}
此处只给出一些初等结果.

\medskip
\phantomsection\label{prop:sga3-VIA-1.1.1}%
\noindent\textbf{命题 1.1.1.}
\textit{设 \(x\) 是一个局部有限型且在 \(A\) 上平坦的
\(A\)-群 \(G\) 的点.
则局部环 \(\mathcal O_{G,x}\) 是 Cohen--Macaulay 环,
且 \(\mathcal O_{G,x}\) 有一个参数系 \(a_1,\ldots,a_n\),
使得 \(\mathcal O_{G,x}/(a_1,\ldots,a_n)\) 是有限平坦
\(A\)-模 (因而有限自由).}

\medskip
先假设 \(A\) 等于其剩余域 \(k\).
此时只需证明 \(\mathcal O_{G,x}\) 是 Cohen--Macaulay 环,
且可以限于 \(x\) 是闭点的情形
(参见 EGA~\(0_{\mathrm{IV}}\), 16.5.13).
由下文\hyperref[lem:sga3-VIA-1.1.2]{引理 1.1.2},
\(G\) 包含一个使 \(\mathcal O_{G,y}\) 为 Cohen--Macaulay 环的闭点 \(y\).
由 SGA~1, I, \S~9, 这等价于说:
对基域 \(k\) 的每个有限扩张 \(K\),
以及 \(\overline G=G\otimes_k K\) 中位于 \(y\) 上方的每个点
\(\overline y\), 局部环 \(\mathcal O_{\overline G,\overline y}\)
都是 Cohen--Macaulay 环.
若 \(K\) 选得充分大, 即 \(K\) 包含 \(k\) 的一个正规扩张,
该正规扩张包含剩余域 \(\kappa(y)\) 与 \(\kappa(x)\),
则 \(\overline y\) 在 \(K\) 上 (严格) 有理,
\(\overline G\) 中位于 \(x\) 上方的每个点 \(\overline x\) 也如此.%
\footnote{\label{note:sga3-VIA-11}编者注: 事实上, 对 \(K\) 的假设蕴含,
对 \(K\) 的每个扩张 \(L\), 每个 \(k\)-态射
\(\kappa(x)\to L\) (分别地, \(\kappa(y)\to L\))
都通过 \(K\) 分解.
因此, \(G\otimes_k K\) 中位于 \(x\) 或 \(y\) 上方的所有点
都以 \(K\) 为剩余域; 即它们在 \(K\) 上 (严格) 有理.}
由于自同构
\[
  r_{\overline x}\circ(r_{\overline y})^{-1}
\]
将 \(\overline y\) 映到 \(\overline x\),
可知 \(\mathcal O_{\overline G,\overline x}\),
从而 \(\mathcal O_{G,x}\) (SGA~1, I, \S~9), 都是 Cohen--Macaulay 环.

现在重新取任意 \(A\).
以上论证适用于 \(k\otimes_A G\),
所以 \(k\otimes_A\mathcal O_{G,x}\) 是 Cohen--Macaulay 环.
若 \(a_1,\ldots,a_n\) 是 \(\mathcal O_{G,x}\) 的一个元素序列,
其在 \(k\otimes_A\mathcal O_{G,x}\) 中的像是一个参数系,
则由 SGA~1, IV, 5.7 或 EGA~\(0_{\mathrm{IV}}\), 15.1.16 可知,
\(a_1,\ldots,a_n\) 是一个 \(\mathcal O_{G,x}\)-正则序列,
且 \(\mathcal O_{G,x}/(a_1,\ldots,a_n)\) 在 \(A\)
上有限且平坦 (因而有限且自由).

\medskip
\phantomsection\label{lem:sga3-VIA-1.1.2}%
\noindent\textbf{引理 1.1.2.}
\textit{每个在 Artin 环 \(A\) 上局部有限型的非空概形 \(X\),
都包含一个局部环为 Cohen--Macaulay 环的闭点 \(x\).}

\medskip
显然可以假设 \(X\) 是仿射的, 其代数为 \(B\),
并对 \(\dim X\) 作归纳 (若 \(X\) 离散, 断言显然成立,
因为此时其所有局部环都是 Artin 环).
由于 \(B\) 在 \(A\) 上有限型, 若 \(\dim B>0\),
则 \(B\) 包含一个既不可逆又不是零因子的元素 \(a\).%
\footnote{\label{note:sga3-VIA-12}编者注: 事实上,
\(B\) 是 Noether 型 Jacobson 环 (参见 \textup{[BAC]}, V, \S~3.4).
若每个不可逆元素都是零因子, 则每个素理想都是 \(B\) 的伴随素理想;
特别地, \(B\) 只有有限个极大理想 \(\mathfrak m_1,\ldots,\mathfrak m_n\).
由于 \(B\) 是 Jacobson 环, 这些 \(\mathfrak m_i\) 的交是 \(B\) 的幂零根.
因此, 每个 \(\mathfrak m_i\) 都是 \(B\) 的极小素理想,
从而 \(\dim B=0\).}
此时 \(X\) 的闭子概形
\[
  X'=\operatorname{Spec} B/(a)
\]
的维数严格小于 \(\dim X\).
由归纳假设, 它包含一个使 \(\mathcal O_{X',x}\)
为 Cohen--Macaulay 环的闭点 \(x\).
由于 \(\mathcal O_{X',x}=\mathcal O_{X,x}/(a)\),
且 \(a\) 在 \(\mathcal O_{X,x}\) 中既不可逆又不是零因子,
环 \(\mathcal O_{X,x}\) 是 Cohen--Macaulay 环
(另见 EGA~IV\textsubscript{2}, 6.11.3).

\medskip
\phantomsection\label{prop:sga3-VIA-1.2}%
\noindent\textbf{命题 1.2.}
\textit{设 \(A\) 是 Artin 局部环,
\(G\) 是一个局部有限型且在 \(A\) 上平坦的 \(A\)-群,
\(x\) 是 \(G\) 的一个闭点.
则存在一个在 \(A\) 上有限且自由的局部 \(A\)-代数 \(A'\),
使得 \(G\otimes_A A'\) 中位于 \(x\) 上方的每个点都在 \(A'\) 上严格有理.}

\medskip
\noindent\footnote{\label{note:sga3-VIA-13}编者注: 以下内容对原文作了扩充.}
事实上, 设 \(k_1\) 是 \(k\) 的一个有限正规扩张,
包含 \(x\) 的剩余域 \(\kappa(x)\).
由\SGARefLink{lem:sga3-V-4.1.2}{引理 V.4.1.2},
存在一个在 \(A\) 上有限且自由、剩余域为 \(k_1\) 的局部 \(A\)-代数 \(A_1\).
此时 (参见编者注~\ref{note:sga3-VIA-11}),
\(G\otimes_A A_1\) 中位于 \(x\in G\) 上方的所有点
\(g_1,\ldots,g_n\) 都以 \(k_1\) 为剩余域;
即 \(g_1,\ldots,g_n\) 在 \(A_1\) 上按第 V 篇 \S~4.e) 的意义是有理的.

设 \(B_1,\ldots,B_n\) 是 \(g_1,\ldots,g_n\) 的局部环.
由 \hyperref[prop:sga3-VIA-1.1.1]{1.1.1},
环 \(B_1,\ldots,B_n\) 有商 \(B'_1,\ldots,B'_n\),
这些商是 Artin 环, 且在 \(A_1\) 上有限且自由. 置
\[
  A'=B'_1\otimes_{A_1}\cdots\otimes_{A_1}B'_n.
\]
则 \(A'\) 是局部环, 在 \(A_1\) 上有限且自由,
并且对每个 \(i=1,\ldots,n\), 有满同态
\phantomsection\label{eq:sga3-VIA-1.2-surjections}%
\[
  B_i\otimes_{A_1}A'
    \twoheadrightarrow B'_i\otimes_{A_1}A'
    \twoheadrightarrow A',
\]
第二个同态由乘法映射
\[
  B'_i\otimes_{A_1}B'_i\twoheadrightarrow B'_i.
\]
诱导. 因此, \(A'\) 具有所需性质.

\medskip
\phantomsection\label{par:sga3-VIA-1.3}%
\noindent\textbf{1.3.}
设 \(e\) 是 \(G\) 的单位元 (或原点),
即 \(\operatorname{Spec}A\) 的唯一点在单位截面
\(\operatorname{Spec}A\to G\) 下的像.
按定义, \(e\) 在 \(A\) 上严格有理.

\par\medskip
\noindent\phantomsection\label{prop:sga3-VIA-1.3.1}%
\textbf{命题 1.3.1.}%
\footnote{\label{note:sga3-VIA-14}编者注: 新增了
\hyperref[concl:sga3-VIA-1.3.1-1]{第 (1) 部分},
它在下文\hyperref[ex:sga3-VIA-1.3.2]{例 1.3.2} 中有用.}
\textit{设 \(G\) 是一个在 Artin 环 \(A\) 上局部有限型且平坦的群,
\(K\) (分别地, \(\overline k\)) 是 \(A\) 的剩余域 \(k\)
的完美闭包 (分别地, 一个代数闭包).}

\begin{enumerate}
\item[(1)]\phantomsection\label{concl:sga3-VIA-1.3.1-1}
\textit{对 \(\overline G=G\otimes_A\overline k\) 的每个闭点 \(x\),
\footnote{Canon 来源说明 (SGA-SRC-000859): 在明确归约为 \(A=k\) 之前,
重版在此处及证明中都将张量积的基印为 \(k\).
这两个对象都定义在 Artin 环 \(A\) 上;
后文归约后的张量积之基仍为 \(k\).}
局部环 \(\mathcal O_{\overline G,e}\) 与 \(\mathcal O_{\overline G,x}\) 同构.
特别地, 切空间 \(T_e\overline G\) 与 \(T_x\overline G\) 的维数相同.}

\item[\SGAInlineRefTarget{prop:sga3-VIA-1.3.1:item:2}(2)] \textit{以下断言等价:}
  \begin{enumerate}
  \item[(i)] \textit{\(G\otimes_A K\) 约化.}
  \item[(i bis)] \textit{\(\mathcal O_{G,e}\otimes_A K\) 约化.}
  \item[\SGAInlineRefTarget{prop:sga3-VIA-1.3.1:item:ii}(ii)] \textit{\(G\) 在 \(A\) 上平滑.}
  \item[(ii bis)] \textit{\(G\) 在原点处于 \(A\) 上平滑.}
  \end{enumerate}
\end{enumerate}

\noindent\textit{证明.}
\SGARefLink{concl:sga3-VIA-1.3.1-1}{(1)}
设 \(x\) 是 \(\overline G\) 的一个闭点.
恰有一个 \(\overline k\)-态射
\[
  s:\operatorname{Spec}\overline k\longrightarrow\overline G
\]
以 \(x\) 为其像.
于是右平移 \(r_s\) 诱导从
\[
  \mathcal O_{\overline G,e}
    =\mathcal O_{G,e}\otimes_A\overline k
\]
到 \(\mathcal O_{\overline G,x}\) 的一个同构,
这就证明了 \SGARefLink{concl:sga3-VIA-1.3.1-1}{(1)}.

现在证明 \SGARefLink{prop:sga3-VIA-1.3.1:item:2}{(2)}.
由 SGA~1, II.2.1, 立即归约到 \(A\) 是域 (\(A=k\)) 的情形.
蕴含关系
\[
  \text{(i)}\Rightarrow\text{(i bis)},\qquad
  \text{(ii)}\Rightarrow\text{(ii bis)},\qquad
  \text{(ii)}\Rightarrow\text{(i)},\qquad
  \text{(ii bis)}\Rightarrow\text{(i bis)}
\]
是显然的, 因此只需证明 (i bis) 蕴含
\SGARefLink{prop:sga3-VIA-1.3.1:item:ii}{(ii)}.

现在, (i bis) 蕴含 \(\mathcal O_{G,e}\otimes_k\overline k\) 约化.
因此, 由 \SGARefLink{concl:sga3-VIA-1.3.1-1}{(1)},
对 \(\overline G\) 的每个闭点 \(x\),
\(\mathcal O_{\overline G,x}\) 都约化,
从而 \(\overline G\) 约化.
由于 \(\overline G\) 在 \(\overline k\) 上局部有限型,
至少存在一个使 \(\mathcal O_{\overline G,y}\) 正则的闭点 \(y\).
由 \SGARefLink{concl:sga3-VIA-1.3.1-1}{(1)},
\(\overline G\) 的所有闭点处的局部环都同构于 \(\mathcal O_{\overline G,e}\);
因此这些局部环全都正则.
所以 \(\overline G\) 在 \(\overline k\) 上平滑,
从而 \(G\) 在 \(k\) 上平滑.

\footnote{\label{note:sga3-VIA-15}编者注: 以下句子及
\hyperref[ex:sga3-VIA-1.3.2]{例 1.3.2} 为新增内容.}
现在可以给出下文由 M. Raynaud 指出的例子:
非完美域 \(k\) 上的群概形 \(G\),
其 \(G_{\mathrm{red}}\) 不是一个 \(k\)-群概形.

% Translation coverage: entire assigned component, retaining all source statements and proofs.
% Statistics: source lines=233; target lines=239; inline math=192; display math=8; unresolved translation TODOs=0; independent compilation not performed here.
% Independent prose review: source and translation compared paragraph by paragraph through this component; statements, hypotheses, proofs, notes and diagram captions covered; bibliography reserved for its assigned reviewer; source math and reference keys unchanged.
